\section{Learned Projection Integration}
\label{sec:projection-integration}

\subsection{What projection means}

Flatten batch and sequence into $M=BS$, let the model hidden width be $K$,
and concatenate the Q/K/V weights:
\begin{equation}
  W_{QKV}=
  \begin{bmatrix}W_Q\\W_K\\W_V\end{bmatrix},\qquad
  Y_{QKV}=XW_{QKV}^{\mathsf T}=[Q\mid K\mid V].
  \label{eq:qkv-projection}
\end{equation}
This is a learned GEMM with reduction dimension $K$.  It is separate from the
attention score contraction in Equation~\ref{eq:score-contraction-summary},
whose reduction dimension is the head dimension 192.

The previous integration ceiling assumed that an upstream projection could
somehow emit the adaptive E2M1 operands.  The current implementation makes
that producer real: a persistent SM100 NVFP4 GEMM computes
Equation~\ref{eq:qkv-projection} and publishes the forward-native Q/K E2M1
payloads, the backward Q/K views, and transposed MXFP4 V directly from its
output register fragments.  BF16 Q/K/V stores are optional rather than part
of the low-precision contract.  Q and K have per-head depth 192; V has depth
128.  Packed weight rows are ordered as all Q heads, all K heads, then all V
heads.

\subsection{NVFP4 learned-projection equation}

For a 16-value activation block $b=(i,c)$, containing one row and columns
$16c{:}16(c+1)$, the standalone quantizer selects one global dequantizer
$g_X$ and one local E4M3 dequantizer $s^X_{i,c}$.  In simplified notation,
\begin{align}
  g_X &\approx \frac{\lVert X\rVert_\infty}{6\cdot448},\\
  s^X_{i,c} &= Q_{\mathrm{E4M3}}\!\left(
    \frac{\max_{16c\le k<16(c+1)}|X_{ik}|}{6g_X}
  \right),\\
  X^4_{ik} &= Q_{\mathrm{E2M1}}\!\left(
    \frac{X_{ik}}{g_Xs^X_{i,c}}
  \right),\qquad
  \widehat X_{ik}=g_Xs^X_{i,c}X^4_{ik}.
  \label{eq:nvfp4-projection-operand}
\end{align}
Activations retain this $1\times16$ rule because each token row is an
independent sample.

Learned weights instead use a transpose-consistent $16\times16$ rule.  For
$B_{r,c}=\{(j,k):16r\le j<16(r+1),\ 16c\le k<16(c+1)\}$,
\begin{align}
  s^W_{r,c}
    &=Q_{\mathrm{E4M3}}\!\left(
      \frac{\max_{(j,k)\in B_{r,c}}|W_{jk}|}{6g_W}
    \right),\\
  W^4_{jk}
    &=Q_{\mathrm{E2M1}}\!\left(
      \frac{W_{jk}}{g_Ws^W_{r,c}}
    \right),\qquad
  \widehat W_{jk}=g_Ws^W_{r,c}W^4_{jk}.
  \label{eq:nvfp4-projection-weight-2d}
\end{align}
The same local scale is repeated into the sixteen hardware $1\times16$
scale records covered by $B_{r,c}$.  Consequently packing $W^{\mathsf T}$
does not choose a second, incompatible set of scales:
\begin{equation}
  \mathcal Q^{2D}_{\mathrm{NV4}}(W^{\mathsf T})
  =\Pi_{\mathsf T}\!\left(
    \mathcal Q^{2D}_{\mathrm{NV4}}(W)
  \right),
  \label{eq:nvfp4-weight-transpose-consistency}
\end{equation}
where $\Pi_{\mathsf T}$ transposes payload blocks and their scale page.  This
is a representation invariant for forward and dgrad, rather than a claim
that the larger block minimizes one-shot quantization error.

The persistent projection computes
\begin{equation}
  \widehat Y_{ij}
  =\sum_{k=1}^{K}\widehat X_{ik}\widehat W_{jk}
  =g_Xg_W
   \sum_b
   \mma_{\mathrm{E2M1}}\!\left(
     X^4_{i,b},W^4_{j,b};s^X_{i,b},s^W_{j,b}
   \right),
  \label{eq:nvfp4-projection}
\end{equation}
accumulates in FP32 TMEM and multiplies by the two global dequantizers.  The
implementation uses M256$\times$N256 output tiles, K256 reduction steps, a
two-CTA cluster, and a four-stage input ring.  Each output fragment is rounded
once to BF16 before its optional BF16 store and low-precision publication, so
all published views share one rounding point.

\subsection{Pair-native RoPE in the projection epilogue}

The first composed forward/backward evaluation omitted rotary position
embedding and therefore did not represent the Llama attention boundary.  The
unified QKV projection now optionally applies full-head RoPE to Q and K before
either BF16 or low-precision publication.  For $D_{QK}=192$, write the usual
split-half rotary coordinates of one unrotated head as
\begin{equation}
  u=[a_0,\ldots,a_{95},b_0,\ldots,b_{95}]^{\mathsf T},
  \qquad
  \Pi u=[a_0,b_0,a_1,b_1,\ldots,a_{95},b_{95}]^{\mathsf T}.
  \label{eq:rope-pair-permutation}
\end{equation}
The Q and K projection-weight rows are converted once to the $\Pi$ order.
Because $\Pi$ is orthogonal and is common to Q and K,
\begin{equation}
  (\Pi q)^{\mathsf T}(\Pi k)=q^{\mathsf T}k,
  \label{eq:rope-permutation-dot-invariance}
\end{equation}
so this storage choice does not alter the unrotated score operator.  It makes
each rotary plane local to one adjacent BF16 pair in the projection epilogue.

For token $t$ and rotary plane $p$, define
\begin{equation}
  R_{t,p}=
  \begin{bmatrix}
    c_{t,p} & -s_{t,p}\\
    s_{t,p} &  c_{t,p}
  \end{bmatrix},
  \qquad
  \widetilde u_{t,p}=R_{t,p}(\Pi u_t)_p.
  \label{eq:pair-native-rope}
\end{equation}
The epilogue applies Equation~\ref{eq:pair-native-rope} to the rounded Q/K
register fragment and then forms every consumer representation from that
rotated fragment:
\begin{align}
  \widetilde Q_R&=\mathcal R_t(\Pi Q_R),
  &C_Q&=Q_{\mathrm{E2M1}}(m^Q_h\widetilde Q_R),\\
  \widetilde K_R&=\mathcal R_t(\Pi K_R),
  &C_K&=Q_{\mathrm{E2M1}}(m^K_h\widetilde K_R).
  \label{eq:projection-rope-then-quantize}
\end{align}
Consequently the optional BF16 Q/K tensors, forward score operands, and all
backward Q/K byte layouts represent the same rotated operator.  V is not
modified.  The packed implementation evaluates both coordinates with one
BF16x2 multiply and one BF16x2 fused multiply-add; this deliberately accepts
one additional BF16 rounding at the already-BF16 publication boundary rather
than expanding every live pair into FP32 registers.

Attention backward returns derivatives with respect to the rotated tensors.
Before projection dgrad or weight-gradient GEMMs, each Q/K pair is mapped
back to the pair-native projection output by
\begin{equation}
  d(\Pi u_t)_p=R_{t,p}^{\mathsf T}d\widetilde u_{t,p}.
  \label{eq:inverse-rope-gradient}
\end{equation}
The reference implementation applies Equation~\ref{eq:inverse-rope-gradient}
in-place to the Q and K fields of the interleaved BF16
$[dQ\mid dK\mid dV]$ tensor; dV is bit-preserved.  A model may keep Q/K
weights and their optimizer state in pair-native row order.  Export to a
split-half checkpoint requires only the inverse row permutation $\Pi^{\mathsf
T}$.

A projection-dgrad specialization now fuses the same inverse map into its
delayed-scale NVFP4 operand producer.  For a 16-value projection-input block
$b$, it publishes
\begin{equation}
  C^{\nabla,4}_{t,b}
  =Q_{\mathrm{E2M1}}\!\left(
    \frac{\mathcal R_t^{\mathsf T}
      [d\widetilde Q_t\mid d\widetilde K_t\mid dV_t]_b}
         {g_\nabla s^\nabla_{t,b}}
  \right),
  \label{eq:inverse-rope-gradient-nvfp4-handoff}
\end{equation}
where $\mathcal R_t^{\mathsf T}$ acts on Q/K pairs and is the identity on V.
The transform occurs after the attention boundary's BF16 rounding and before
local amax selection, so its E2M1 payload and scale page are bit-exact with
the reference sequence ``inverse RoPE, then NVFP4 quantize.''  An optional
store republishes the inverse-rotated BF16 tile for projection weight-gradient
consumers.  That store preserves the reference bits and is not a performance
penalty in the measured kernel: its asynchronous BF16 publication overlaps
the packing tail.  Both forms are retained, with republication selected when
projection weight-gradient consumers require the inverse-rotated BF16 tile.

\subsection{One Q/K payload shared by forward and backward}

For a BF16-rounded Q register slice $Q_R$ and its head multiplier $m^Q_h$, the
epilogue computes
\begin{equation}
  C_Q=Q_{\mathrm{E2M1}}(m^Q_hQ_R).
  \label{eq:projection-epilogue-q}
\end{equation}
The same codes are rearranged into:
\begin{align}
  \mathcal L_{Q,\mathrm{seq}}(C_Q)
    &\in\mathrm{uint8}^{B\times H\times192\times S},\\
  \mathcal L_{Q,\mathrm{depth}}(C_Q)
    &\in\mathrm{uint8}^{B\times H\times S\times96}.
  \label{eq:q-layout-functions}
\end{align}
For K it independently forms $C_K=Q_{\mathrm{E2M1}}(m^K_hK_R)$ and publishes
\begin{align}
  \mathcal L_{K,\mathrm{aligned}}(C_K)
    &\in\mathrm{uint8}^{B\times S\times H\times192},\\
  \mathcal L_{K,\mathrm{depth}}(C_K)
    &\in\mathrm{uint8}^{B\times H\times S\times96}.
  \label{eq:k-layout-functions}
\end{align}
The compact tensors $\mathcal L_{Q,\mathrm{depth}}$ and
$\mathcal L_{K,\mathrm{depth}}$ are the storage allocation, not copies made
for backward.  Forward obtains typed E2M1 views of the same bytes while
backward consumes uint8 views.  The forward scale hierarchy is
\begin{align}
  s^{Q}_{\mathrm{local}}=s^{K}_{\mathrm{local}}&=448,\\
  g^Q_h&=\frac{1}{448m^Q_h},\qquad
  g^K_h=\frac{1}{448m^K_h},\\
  \widehat Q_R=g^Q_hs^Q_{\mathrm{local}}C_Q&=\frac{C_Q}{m^Q_h},\qquad
  \widehat K_R=g^K_hs^K_{\mathrm{local}}C_K=\frac{C_K}{m^K_h}.
  \label{eq:shared-qk-forward-reconstruction}
\end{align}
Thus forward and backward consume one represented Q/K operator.  There is no
second forward payload and no conversion between the two APIs.

Together with the optional BF16 Q/K store, one register slice can therefore
emit the ordinary tensor and its consumer-specific E2M1 views.
Aligned-U4 rows use eight E2M1 bytes followed by an eight-byte hole for every
16 logical values.  Compact depth rows pack two E2M1 values per byte without
those holes.

The four low-precision tensors are bit-exact against the standalone
precomputed-scale packer.  Fused and separate publication therefore have
identical accuracy; their difference is only movement and scheduling.

\subsection{Single-quantization pure-FP4 Q/K fanout}

The pure-FP4 backward needs two additional compact Q/K operand views for its
block-scaled dQ and dK instructions.  A naive unified producer first formed
the four adaptive forward/backward views above, then quantized the same Q/K
register fragment a second time at the pure route's fixed multiplier.  The
aggressive specialization instead performs one fixed-scale quantization,
\begin{equation}
  C_X^{(16)}=Q_{\mathrm{E2M1}}(16X_R),\qquad X\in\{Q,K\},
  \label{eq:pure-projection-single-quant}
\end{equation}
and fans those codes directly into every consumer layout,
\begin{equation}
  \mathcal B_{QK}^{\mathrm{pure}}
  =\left\{\Pi_\ell\!\left(C_{\chi(\ell)}^{(16)}\right)\right\}_{\ell=1}^{6},
  \qquad \chi(\ell)\in\{Q,K\}.
  \label{eq:pure-projection-layout-fanout}
\end{equation}
Here the six members are the four aligned/compact score and gradient layouts
shared with the ordinary API plus the compact dQ/dK matrix operands.  The
layout maps $\Pi_\ell$ change byte placement only; no member requantizes the
fragment.  The fixed representation also removes the adaptive seven-word
head record from this specialization.  Its forward scale hierarchy is
\begin{equation}
  s^Q_{\mathrm{local}}=s^K_{\mathrm{local}}=448,
  \qquad g^Q=g^K=\frac{1}{448\cdot16},
  \label{eq:pure-projection-forward-scales}
\end{equation}
while the compact backward dQ/dK operands need no materialized scale tensors.
All six payloads and the forward scale pages are bit-exact against the
standalone fixed-scale publisher.  This is an intentionally pure-only
contract: adaptive FP4+FP8 continues to use the robust per-head record.

\subsection{Projection-native MXFP4 V}

For the BF16-rounded projection fragment, define the two-dimensional tile
$T_{u,v}=\{(t,d):32u\le t<32(u+1),\ 32v\le d<32(v+1)\}$.  The retained
MXFP4 policy forms one E8M0 scale for the complete tile,
\begin{align}
  s^V_{u,v}&=Q_{\mathrm{E8M0}}\!\left(
    \frac{\max_{(t,d)\in T_{u,v}}|V_{R,t,d}|}{6}
  \right),\\
  C^V_{t,d}&=Q_{\mathrm{E2M1}}\!\left(
    \frac{V_{R,t,d}}{s^V_{u,v}}
  \right),\qquad
  \widehat V_{t,d}=s^V_{u,v}C^V_{t,d}.
  \label{eq:projection-native-v}
\end{align}
The scale is repeated into the thirty-two hardware $1\times32$ scale records
inside $T_{u,v}$.  The same codes and scale are then written in both the
feature-major layout required by forward PV and the row-major layout required
by backward dP.  Therefore
\begin{equation}
  \mathcal B^{V}_{\mathrm{bwd}}
  =\Pi_{SD}\!\left(\mathcal B^{V}_{\mathrm{fwd}}\right),
  \label{eq:mxfp4-v-transpose-consistency}
\end{equation}
with no independent row-scale selection after the transpose.  This removes
the BF16 V reread and standalone row-to-column packing from forward setup and
makes the two MXFP4 consumers observe one represented V operator.  The legacy
$1\times32$ policy remains available only as an explicit diagnostic.

The complete optional publication contract is
\begin{equation}
  \mathcal B_{QKV}=
  \left\{
    C_Q,C_K,s^Q,g^Q,s^K,g^K,C^V,s^V,
    \mathcal L_{Q,\mathrm{seq}},\mathcal L_{K,\mathrm{aligned}},
    m^Q_h,m^K_h
  \right\}
  \cup
  \left\{Q_{\mathrm{BF16}},K_{\mathrm{BF16}},V_{\mathrm{BF16}}\right\}_{\mathrm{optional}}.
  \label{eq:unified-qkv-bundle}
\end{equation}
When every downstream consumer accepts $\mathcal B_{QKV}$, disabling the
three BF16 stores avoids $BSH(192+192+128)\cdot2=BSH\cdot1024$ bytes of
global publication.

\subsection{Scale ownership at the composed D64 boundary}

The projection epilogue also publishes E4M3 Q, K, and V for attention
backward.  These are deliberately independent of the forward FP4 multipliers
$m^Q_h$ and $m^K_h$.  Let $\lambda=2^2=4$ be the fixed E4M3 publication
factor and let $G=dO$ denote the loss-scaled upstream gradient.  Backward
consumes
\begin{equation}
  Q'=\lambda Q,\qquad K'=\lambda K,\qquad
  V'=\lambda V,\qquad G'=\lambda G,
  \label{eq:d64-fixed-fp8-publication}
\end{equation}
and compensates the score product with
\begin{equation}
  \alpha'=\frac{\alpha}{\lambda^2}=\frac{\alpha}{16}.
  \label{eq:d64-fixed-fp8-score-scale}
\end{equation}
The represented softmax is consequently unchanged.  Applying the chain rule
in encoded units gives
\begin{equation}
  dV'=\lambda dV,\qquad dP'=\lambda^2dP,\qquad
  dQ'=\lambda dQ,\qquad dK'=\lambda dK.
  \label{eq:d64-fixed-fp8-gradient-units}
\end{equation}
If dynamic loss scaling contributes $L$, inverse RoPE and the QKV projection
consume each attention-gradient field with $1/(\lambda L)$.  Optional
calibration gains $c_Q,c_K,c_V$ remain field-local,
\begin{equation}
  \bar dQ=\frac{c_Q}{\lambda L}R^{\mathsf T}dQ',\qquad
  \bar dK=\frac{c_K}{\lambda L}R^{\mathsf T}dK',\qquad
  \bar dV=\frac{c_V}{\lambda L}dV'.
  \label{eq:d64-field-local-gradient-decode}
\end{equation}
The default is $c_Q=c_K=c_V=1$.  A forward probability approximation does
not implicitly own these gains: any non-unit correction must be exported by
the topology and validated per field.  In particular, a historical common
0.632 multiplier is not inferred from the softmax mode because it suppresses
Q/K/V gradient norms without changing the output-projection weight gradient.

\subsection{Adaptive scale production remains a contract}

Outside the fixed-scale pure specialization, the fused epilogue
\emph{consumes} the seven-word per-head scale record; it does not calculate
Equation~\ref{eq:head-stats} from its distributed output tiles.  A true
same-launch amax/RMS reduction would need cross-CTA completion before earlier
Q/K tiles could be finalized, defeating the simple streaming epilogue.

The current API consequently exposes three legal scale sources:

\begin{enumerate}
  \item the standalone robust Q/K reduction, used for validation;
  \item precomputed device-resident scales from an upstream producer or
  calibration policy;
  \item a future learned-projection epilogue with a delayed, pipelined, or
  model-defined scale contract.
\end{enumerate}

Current projection validation derives the record from BF16 reference Q/K
before timing the projection hot stage.  The measured epilogue is real, but
this detail prevents interpreting its timing as a solved same-step adaptive
scale reduction.  An integrated model must state how $m^Q_h$ and $m^K_h$ are
supplied.

\subsection{Upstream low-precision dO contract}

The backward also needs two MXFP4 orientations of the output gradient
$G=dO$ and the softmax centering term.  For row blocks $b_r$ and transposed
column blocks $b_c$, the producer emits
\begin{align}
  C^{G,r}_{i,b_r}&=Q_{\mathrm{E2M1}}
    \!\left(G_{i,b_r}/s^{G,r}_{i,b_r}\right),\\
  C^{G,c}_{d,b_c}&=Q_{\mathrm{E2M1}}
    \!\left(G^{\mathsf T}_{d,b_c}/s^{G,c}_{d,b_c}\right),\\
  \delta_i&=\sum_{d=1}^{128}O_{i,d}G_{i,d}.
  \label{eq:upstream-dout-contract}
\end{align}
One fused producer reads dO once and publishes both code/scale orientations
and $\delta$.  This is the upstream boundary used by the all-MX experiments;
the separate row quantizer, column quantizer, and delta kernel remain only as
a bit-exact diagnostic reference.

\subsection{Probability replay boundary}

Within one backward iteration, the retained route reuses the represented
E4M3 probability produced for dV when forming dS rather than reevaluating the
exponential.  If $C^P$ denotes that code,
\begin{equation}
  \widehat{dS}_{ij}=\widehat P_{ij}
    \left(\widehat{dP}_{ij}-\delta_i\right),\qquad
  \widehat P_{ij}=s^P_{ij}C^P_{ij}.
  \label{eq:represented-p-replay}
\end{equation}
This is a local register/TMEM reuse contract, not an HBM cache.  Caching a
complete MXFP4 probability matrix across forward and backward would require
\begin{equation}
  M_P=\frac{BHS^2}{2}
      +BH\left(\frac{S}{128}\right)^2(32\cdot16)
      \quad\text{bytes}
  \label{eq:probability-cache-bytes}
\end{equation}
for payload plus scale pages.  At S8192/H8 this is approximately 272 MiB, so
the write/read traffic overwhelms the measured local replay saving.

\subsection{Reduction-width dispatch}

Publication cost is controlled by projection reduction depth $K$, not by head
count in isolation.  With only four K256 steps at K1024, one consumer
warpgroup cannot hide BF16 rounding, E2M1 conversion, transpose, and four
global layouts behind the tensor producer.  The separate 32-register packer
spreads that work over a wider CTA grid.  With twelve or more K256 steps,
publication hides under the MMA pipeline and avoiding the BF16 HBM reread is
decisive.

The legacy Q/K-only public dispatcher therefore selects
\begin{equation}
  \operatorname{publish}(K)=
  \begin{cases}
    \text{NVFP4 projection + parallel packer},&K<3072,\\
    \text{fused register-epilogue publication},&K\ge3072.
  \end{cases}
  \label{eq:projection-dispatch}
\end{equation}
Both paths return one \code{B300AdaptiveLowpOperands} object and identical
layout bits.

The unified QKV entry point always uses the fused projection epilogue because
V publication and optional BF16 suppression are part of the single producer
contract.  Its projection scratch aliases the Q/K code staging region with
the V BF16-pair staging region: they are disjoint in lifetime.  This reduces
static scratch from 10,704 to 8,592 bytes without changing the hot backward
kernel.

\subsection{Projection backward}

For the PyTorch weight convention in Equation~\ref{eq:learned-projections-summary},
the input-gradient contribution is
\begin{equation}
  dX=dQW_Q+dKW_K+dVW_V.
  \label{eq:projection-dgrad}
\end{equation}
The first integration writes fully reduced BF16 gradients into one per-head
buffer ordered as $[dQ_{192},dK_{192},dV_{128}]$.  A persistent weight packing
step stores $[W_Q,W_K,W_V]$ in the matching order.  One, two, or four
head-contiguous GEMMs consume the buffer, selected by shape, without three
standalone gradient tensors or an explicit concatenation.

The newer dQ path joins readiness to the projection reduction loop.  Let
$r$ index a 256-row projection tile, $h$ a head, and $o$ a causal K/V-owner
cluster.  Each of the two CTAs in owner $o$ publishes one release increment
only after all four reducer warps have retired their BF16 TMA reductions and
ordered the async proxy:
\begin{align}
  a_{h,t} &\leftarrow a_{h,t}+1,\qquad
  \mathrm{ready}(h,t,r)\iff a_{h,t}\ge 2(r+1),\\
  t&\in\{2r,2r+1\}.
  \label{eq:dq-head-arrival}
\end{align}
The factor $2(r+1)$ is the exact number of contributing CTAs for either
128-row half of causal row pair $r$.  A device-scope acquire poll is placed
immediately before the first K64 slice of each head, not before the complete
projection tile.  The projection accumulator therefore evaluates
\begin{equation}
  dX_{R,C}=\sum_{h=0}^{H-1}\sum_{j=0}^{2}
    dQ_{R,h,64j:64(j+1)}
    W_{Q,C,h,64j:64(j+1)}^{\mathsf T},
  \label{eq:head-pipelined-dq-projection}
\end{equation}
waiting for head $h$ only when its first $j=0$ slice reaches the load
producer.  This is a reduction-topology integration: readiness is part of the
K loop, rather than a whole-row gate in front of GEMM.

Two FP32 TMEM accumulators alternate output tiles.  While the epilogue drains
slot $p$ to BF16, the tensor warpgroup issues the next output into slot
$p\oplus1$.  This removes the former MMA--epilogue serialization without
allocating TMEM inside the attention kernel; the projection kernel owns all
512 of its own TMEM columns.  The consumer remains at 80 registers with zero
spills, and the attention specialization remains at 128 registers with zero
spills.

A bounded early grid is necessary so projection does not displace the
attention owner grid.  For wide outputs, that bounded grid cannot expand when
attention retires, so an experimental split computes a ready row prefix with
the persistent consumer and a disjoint suffix with the vendor GEMM:
\begin{equation}
  dX=\begin{bmatrix}
    \mathcal P_{\theta}(dQ_{0:R},W_Q)\\
    \mathrm{GEMM}(dQ_{R:S},W_Q)
  \end{bmatrix}.
  \label{eq:dq-projection-tail-split}
\end{equation}
The two streams write disjoint rows.  At S8192/H8 the all-persistent route now
slightly beats attention plus cuBLAS.  H16/H24/H64 still favor the ordinary
BF16 reduction followed by the vendor GEMM, so the public entry point
shape-dispatches to that non-regressing fallback unless the pipelined route is
explicitly forced for profiling.

\subsection{Absorbing a hierarchical dQ sum}

The next experiment removes the \emph{completed} standalone dQ tensor from
the projection contract.  Partition K/V owners by parity and define two BF16
reduction lanes
\begin{equation}
  dQ^{(\ell)}=\sum_{o:\,o\bmod2=\ell}dQ^{(o)},\qquad
  dQ=dQ^{(0)}+dQ^{(1)}.
  \label{eq:dq-two-lane-hierarchy}
\end{equation}
The attention reducer publishes each owner into its compile-time-selected
lane.  Projection then exploits linearity:
\begin{align}
  dX_Q
    &=\left(dQ^{(0)}+dQ^{(1)}\right)W_Q\\
    &=dQ^{(0)}W_Q+dQ^{(1)}W_Q.
  \label{eq:dq-hierarchy-projection}
\end{align}
Both lane products accumulate into the same FP32 TMEM output.  No kernel ever
constructs or consumes a global BF16 array equal to $dQ$.  The ordinary
one-lane producer is a separate template instantiation, so hierarchy indexing
does not add a modulo or branch to the retained path.

The construction is numerically sound but does not improve time.  At
S8192/H8, attention plus cuBLAS measured 0.773824~ms and the two-lane consumer
0.805888~ms, a 32.064~$\mu$s loss, with dX cosine 0.99999422.  At S4096/H24,
the corresponding times were 0.736000~ms and 1.063392~ms, a
327.392~$\mu$s loss, with cosine 0.99999648.  Both projection specializations
compile at 80 registers, two barriers, and zero spills.

Equation~\ref{eq:dq-hierarchy-projection} duplicates the complete weight MMA.
A second implementation instead formed $dQ^{(0)}+dQ^{(1)}$ in shared memory
before one MMA.  Scalar and packed-BF16 serial folds took approximately 2.68
and 2.33~ms.  Moving the fold to idle producer warps and using explicit
128-bit shared transfers improved it to 0.860608~ms at S8192/H8, but remained
slower than the two-MMA form.  The tensor-native hierarchy is therefore kept
only as an opt-in structural diagnostic.

This distinction matters: the experiment eliminates a completed BF16 dQ
\emph{tensor}, but not BF16 global partial traffic.  The producer still emits
the same number of values, split across two arrays, and projection pays either
an additional weight product or an explicit fold.  A winning hierarchy must
also change the pre-publication representation:
\begin{equation}
  \text{owner partials}
  \xrightarrow{\text{on-chip grouping or compact FP8/FP4}}
  \text{projection reduction}.
  \label{eq:dq-winning-hierarchy-contract}
\end{equation}
A winning path performs Equation~\ref{eq:dq-winning-hierarchy-contract}
without a BF16 dQ round trip.  An on-chip super-owner can reduce several K/V
owners before one global write.
Alternatively, compact low-precision partials can make repeated projection
products proportional to the higher FP8/FP4 tensor rate.  Reversing the whole
kernel to q-centric ownership completes dQ locally but transfers the global
reduction to the combined 320 dK/dV features rather than the 192 dQ features,
so it is not the preferred next topology.

\subsection{Stacked low-precision QKV-gradient handoff}

The final experiment combines the three projection-gradient fields before one
NVFP4 projection.  Let $R_t^{\mathsf T}$ apply inverse RoPE to one Q or K
head at token $t$.  For the two-lane diagnostic, the row presented to the
quantizer is
\begin{equation}
  G_t^{(2)}=
  \left[
    R_t^{\mathsf T}\!\left(dQ_t^{(0)}+dQ_t^{(1)}\right)
    \;\middle|\;
    R_t^{\mathsf T}dK_t
    \;\middle|\;
    dV_t
  \right].
  \label{eq:stacked-qkv-gradient-row}
\end{equation}
The producer loads a K128 tile, assigns its two K64 halves to independent row
groups, performs the dQ lane addition and inverse rotations in registers, and
chooses one delayed-scale NVFP4 representation per K16 block:
\begin{align}
  s^G_{t,b}
    &=Q_{\mathrm{E4M3}}\!\left(
      \frac{\max_{k\in b}|G^{(2)}_{t,k}|}{6\bar g_G}
    \right),\\
  C^G_{t,k}
    &=Q_{\mathrm{E2M1}}\!\left(
      \frac{G^{(2)}_{t,k}}{\bar g_Gs^G_{t,b}}
    \right),\qquad
  \widehat G_{t,k}=\bar g_Gs^G_{t,b}C^G_{t,k}.
  \label{eq:stacked-qkv-gradient-nvfp4}
\end{align}
With the correspondingly stacked persistent weight
$W_\Sigma=[W_Q^{\mathsf T},W_K^{\mathsf T},W_V^{\mathsf T}]^{\mathsf T}$,
one projection computes
\begin{equation}
  \widehat{dX}_{t,j}
  =\bar g_Gg_W\sum_b
    \mma_{\mathrm{E2M1}}\!\left(
      C^G_{t,b},C^W_{j,b};s^G_{t,b},s^W_{j,b}
    \right).
  \label{eq:stacked-qkv-gradient-projection}
\end{equation}
This form avoids a completed dQ array, but enabling per-tile readiness changes
the attention schedule.  Even a one-lane signaling control slows the attention
kernel.  The retained activation-only variant therefore uses the ordinary
unsignaled BF16 dQ endpoint and replaces $dQ^{(0)}+dQ^{(1)}$ in
Equation~\ref{eq:stacked-qkv-gradient-row} by completed $dQ$.  It still fuses
inverse RoPE, three-field layout conversion, local scaling, and packing before
the single NVFP4 projection.  It does \emph{not} remove BF16 dQ
materialization, and it does not provide the BF16 QKV tensor required by the
projection weight-gradient GEMM.  The two contracts are therefore reported
separately rather than treating the activation-only result as a complete
training step.

\subsection{Compact NVFP4 dQ handoff}

A second implementation changes the representation at the final
attention--projection boundary without changing the owner reduction.  Let
$D=dQ_{\mathrm{complete}}$ be the private BF16 reduction scratch and let
$\bar g_D$ be a device-resident NVFP4 decode scale supplied by calibration or
a previous step.  For each K16 block $b$, the producer forms
\begin{align}
  s^D_{i,b}
    &=Q_{\mathrm{E4M3}}\!\left(
      \frac{\max_{k\in b}|D_{ik}|}{6\bar g_D}
    \right),\\
  C^D_{ik}
    &=Q_{\mathrm{E2M1}}\!\left(
      \frac{D_{ik}}{\bar g_Ds^D_{i,b}}
    \right),
  \qquad
  \widehat D_{ik}=\bar g_Ds^D_{i,b}C^D_{ik}.
  \label{eq:compact-dq-nvfp4}
\end{align}
The matrix-wide amax pass is absent from this hot path.  With a persistently
packed projection weight
$\widehat W_{jk}=g_Ws^W_{j,b}C^W_{jk}$, the consumer evaluates
\begin{equation}
  \widehat{dX}_{Q,ij}
  =\bar g_Dg_W\sum_b
    \mma_{\mathrm{E2M1}}\!\left(
      C^D_{i,b},C^W_{j,b};s^D_{i,b},s^W_{j,b}
    \right).
  \label{eq:compact-dq-nvfp4-projection}
\end{equation}
The public result contains $\widehat{dX}_Q$, dK, and dV; standalone dQ is not
returned.  The generic NVFP4 consumer uses M256$\times$N256 output tiles,
K256 reduction steps, a two-CTA cluster, FP32 TMEM accumulation, and one BF16
output rounding.  It compiles at 86 registers with two barriers and no spills.

This is a compact \emph{final-tile} handoff, not yet a compact owner-partial
reduction.  The unclustered attention owners still require an additive BF16
TMA target, so $D$ exists as private scratch until the last owner completes:
\begin{equation}
  \{dQ^{(o)}\}_o
  \xrightarrow{\mathrm{BF16\ TMA\ reduction}}
  D
  \xrightarrow{Q_{\mathrm{NV4}}}
  (C^D,s^D,\bar g_D)
  \xrightarrow{\mathrm{NV4\ projection}}
  \widehat{dX}_Q.
  \label{eq:compact-dq-current-boundary}
\end{equation}
It nevertheless measures whether compact projection consumption can repay its
conversion.  Delayed scale overestimates from $1\times$ through $8\times$ the
same-step ideal preserve approximately 0.99098 cosine because the local K16
scales renormalize the blocks.  Underestimates clip: $0.5\times$, $0.25\times$,
and $0.125\times$ give approximately 0.9905, 0.9849, and 0.9637 cosine.
Consequently a conservative delayed policy is viable; a winning fused owner
scheme must still publish tile-ready codes or quantize inside a super-owner.

\subsection{The new ceiling}

Unified projection publication removes three previously exposed costs:
\begin{equation}
  \text{BF16 Q/K/V HBM publication}
  \rightarrow
  \text{Q/K/V HBM reread}
  \rightarrow
  \text{standalone pack/transpose/quantize}.
  \label{eq:removed-qk-roundtrip}
\end{equation}
It does not remove the global dQ reduction.  A timing-only attention run that
suppressed dQ publication reached approximately 0.693 ms at S8192/H8, versus
0.746 ms for the direct-BF16 endpoint.  The head-pipelined route now
participates in readiness and reduction scheduling, but still reads the
globally materialized BF16 sum.  Reaching the remaining floor requires
projection backward to absorb the hierarchical sum before that publication,
not merely observe it earlier.  The two-lane experiment shows more precisely
that absorbing only the \emph{last addition} is also insufficient: owner
grouping or compact publication must remove the BF16 producer traffic at the
same time.
